\problema{Longest Palindromic Substring}{5}{src/07-dp/5-longest-palindromic-substring.cpp}{M}

Dada una cadena \texttt{s}, encuentra el substring (secuencia de caracteres
\emph{contiguos}) más largo que sea un palíndromo, es decir, que se lea
igual de izquierda a derecha que de derecha a izquierda.

\paragraph{La idea, con un ejemplo de la vida real.} Un palíndromo siempre
es simétrico respecto a un punto central, como un espejo. Ese centro puede
caer exactamente sobre un carácter (palíndromos de longitud impar, como
\texttt{"aba"}) o justo entre dos caracteres (palíndromos de longitud par,
como \texttt{"abba"}). La técnica de \textbf{expansión desde el centro} (una
variante de \emph{two pointers}, dos punteros) consiste en probar, uno por
uno, cada posible centro de la cadena, y desde ahí ir separando dos punteros
hacia afuera mientras los caracteres que van quedando en los extremos sigan
coincidiendo, igual que cuando revisas si una palabra es capicúa comparando
letras desde el centro hacia las puntas.

\begin{center}
\ttfamily
\begin{tabular}{l}
s = "babad"\\
centro en la 'a' del medio (índice 1) -> se expande a "bab"\\
también existe "aba" como respuesta válida, de la misma longitud\\
\end{tabular}
\end{center}

\paragraph{Cómo lo resuelve el código, paso a paso.}
\begin{enumerate}
  \item Para una cadena de longitud \texttt{n} hay \texttt{2n - 1} centros
        posibles: \texttt{n} centros de un solo carácter y \texttt{n - 1}
        centros ``entre pares'' de caracteres consecutivos.
  \item Para cada centro, colocamos dos punteros \texttt{L} y \texttt{R}
        sobre ese centro (iguales si el centro es un carácter, o adyacentes
        si el centro es ``entre pares'').
  \item Mientras \texttt{L} y \texttt{R} sigan dentro del rango de la cadena
        y \texttt{s[L] == s[R]}, expandimos: \texttt{L} retrocede uno y
        \texttt{R} avanza uno.
  \item Al detenerse la expansión, el palíndromo válido más reciente es el
        substring entre \texttt{L + 1} y \texttt{R - 1}.
  \item Si esa longitud es mayor que la mejor encontrada hasta ahora,
        actualizamos el inicio y la longitud del mejor palíndromo.
  \item Al terminar de probar todos los centros, devolvemos el substring
        guardado como el mejor.
\end{enumerate}

\lstinputlisting{src/07-dp/5-longest-palindromic-substring.cpp}

\complejidad{$O(n^2)$}{$O(1)$ extra}

\paragraph{¿Qué significa esa complejidad?} Hay \texttt{2n - 1} centros
posibles, y cada expansión puede llegar a recorrer, en el peor caso, hasta
\texttt{O(n)} pasos (por ejemplo, cuando toda la cadena es un solo
palíndromo), así que el trabajo total es cuadrático. El espacio extra es
constante porque solo se usan un puñado de variables (los dos punteros y las
marcas del mejor resultado encontrado), sin ninguna tabla auxiliar, a
diferencia de la solución alternativa de DP con tabla 2D, que sí gastaría
espacio cuadrático.

\paragraph{Consejos para la entrevista (Amazon).} Muchos candidatos proponen
primero la tabla 2D de DP (\texttt{dp[i][j]} = ``¿el substring entre
\texttt{i} y \texttt{j} es palíndromo?''), que es correcta pero usa espacio
\texttt{O(n\textsuperscript{2})}; vale la pena mencionarla y luego señalar
que expansión desde el centro llega al mismo tiempo con espacio constante.
Un error común es olvidar los centros ``entre pares'' y solo probar centros
de un carácter, lo que hace fallar casos de palíndromos de longitud par como
\texttt{"cbbd"}. También conviene aclarar qué hacer ante empates de longitud
(cualquier palíndromo de la longitud máxima es una respuesta válida, no hay
una única correcta).
